\documentclass[11pt, a4paper]{article}
\usepackage[a4paper, top=2.5cm, bottom=2.5cm, left=2cm, right=2cm]{geometry}
\usepackage{fontspec}
\usepackage[english, bidi=basic, provide=*]{babel}
\babelprovide[import, onchar=ids fonts]{english}
\babelfont{rm}{TeX Gyre Termes}
\babelfont{sf}{TeX Gyre Heros}
\babelfont{tt}{TeX Gyre Cursor}
\usepackage{enumitem}
\usepackage{amsmath}

\title{\textbf{Chapter 2: The Chemical Context of Life \\ Exam Notes}}
\author{}
\date{}

\begin{document}

\maketitle

\section*{2.1 Matter: Elements and Compounds}
Organisms are composed of matter (anything that takes up space and has mass).
\begin{itemize}
    \item \textbf{Elements:} Substances that cannot be broken down further by chemical reactions. There are 92 naturally occurring elements (e.g., Gold, Carbon, Oxygen).
    \item \textbf{Compounds:} Substances consisting of two or more different elements combined in a fixed ratio (e.g., NaCl, $H_2O$).
    \item \textbf{Emergent Properties:} A compound has characteristics different from those of its constituent elements (e.g., Na is a reactive metal, Cl is a poisonous gas, but NaCl is edible table salt).
\end{itemize}

\subsection*{The Elements of Life}
\begin{itemize}
    \item \textbf{Essential Elements:} Organisms need about 20--25\% of the 92 natural elements to live and reproduce. 
    \item \textbf{The Big Four:} \textbf{Oxygen (O), Carbon (C), Hydrogen (H), and Nitrogen (N)} make up $\sim$96\% of living matter.
    \item \textbf{Trace Elements:} Required by organisms in only minute quantities.
    \begin{itemize}
        \item \textit{Iron (Fe):} Required by all forms of life.
        \item \textit{Iodine (I):} Essential for thyroid hormone in vertebrates; deficiency causes goiter.
    \end{itemize}
    \item \textbf{Toxic Elements:} Some elements are naturally toxic, such as Arsenic. Some plant species (e.g., in serpentine soil) have evolved adaptations to survive in environments with toxic elements like chromium, nickel, and cobalt.
\end{itemize}

\section*{2.2 Atomic Structure and Element Properties}
An \textbf{atom} is the smallest unit of matter that still retains the properties of an element.

\subsection*{Subatomic Particles}
\begin{itemize}
    \item \textbf{Protons ($p^+$):} Positively charged, found in the nucleus. Determine the element's identity.
    \item \textbf{Neutrons ($n^0$):} Electrically neutral, found in the nucleus. Determine the isotope.
    \item \textbf{Electrons ($e^-$):} Negatively charged, form a "cloud" around the nucleus. Determine chemical behavior.
    \item \textbf{Dalton (amu):} Unit of measurement for atomic mass. Protons and neutrons each have a mass of $\sim$1 dalton. Electron mass is negligible.
\end{itemize}

\subsection*{Atomic Number and Mass}
\begin{itemize}
    \item \textbf{Atomic Number:} The number of protons in an atom's nucleus. In a neutral atom, protons = electrons.
    \item \textbf{Mass Number:} The sum of protons + neutrons in the nucleus. 
    \item \textit{Calculation:} Number of neutrons = Mass number $-$ Atomic number.
\end{itemize}

\subsection*{Isotopes and Radiometric Dating}
\begin{itemize}
    \item \textbf{Isotopes:} Different atomic forms of the same element that have the same number of protons but \textbf{different numbers of neutrons} (hence different mass).
    \item \textbf{Radioactive Isotopes:} Unstable isotopes where the nucleus decays spontaneously, giving off particles and energy (e.g., $^{14}C$ decaying into $^{14}N$).
    \item \textbf{Diagnostic Uses:} Used as tracers in metabolism tracking and in PET scans to detect cancerous tissues.
    \item \textbf{Radiometric Dating:} Uses the constant \textbf{half-life} (time taken for 50\% of parent isotope to decay) to determine the age of fossils and rocks.
\end{itemize}

\subsection*{Energy Levels of Electrons}
\begin{itemize}
    \item \textbf{Potential Energy:} Energy that matter possesses due to its location or structure. 
    \item \textbf{Electron Shells:} Electrons exist at fixed levels of potential energy. The further an electron is from the nucleus, the \textbf{greater} its potential energy.
    \item Electrons move between shells by absorbing energy (moving outward) or losing energy (falling inward, often releasing light/heat).
\end{itemize}

\subsection*{Electron Distribution and Orbitals}
\begin{itemize}
    \item \textbf{Valence Electrons:} Electrons in the outermost shell (\textbf{valence shell}). They govern chemical behavior. Atoms with completed valence shells (e.g., Helium, Neon, Argon) are \textbf{inert} (unreactive).
    \item \textbf{Orbitals:} 3D spaces where an electron is found 90\% of the time. Max 2 electrons per orbital.
    \begin{itemize}
        \item 1st shell: One spherical $s$ orbital ($1s$). Holds 2 electrons.
        \item 2nd shell: One $s$ orbital ($2s$) and three dumbbell-shaped $p$ orbitals ($2p$). Holds up to 8 electrons.
    \end{itemize}
\end{itemize}

\section*{2.3 Chemical Bonding: Molecules and Ionic Compounds}
Atoms with incomplete valence shells interact to share or transfer electrons, forming \textbf{chemical bonds}.

\subsection*{Covalent Bonds}
\begin{itemize}
    \item \textbf{Covalent Bond:} The \textbf{sharing} of a pair of valence electrons by two atoms. 
    \item \textbf{Molecule:} Two or more atoms held together by covalent bonds (e.g., $H_2$).
    \item \textbf{Single vs. Double Bonds:} Single bonds share one pair of electrons ($H-H$); double bonds share two pairs ($O=O$).
    \item \textbf{Bonding Capacity (Valence):} Usually equals the number of electrons required to complete the valence shell. ($H=1$, $O=2$, $N=3$, $C=4$).
\end{itemize}

\subsection*{Electronegativity and Polarity}
\begin{itemize}
    \item \textbf{Electronegativity:} An atom's attraction for shared electrons.
    \item \textbf{Nonpolar Covalent Bond:} Electrons are shared equally because atoms have similar electronegativity (e.g., $O_2$, $CH_4$).
    \item \textbf{Polar Covalent Bond:} One atom is more electronegative, pulling shared electrons closer. Creates partial charges ($\delta+$ and $\delta-$) (e.g., $H_2O$, where Oxygen is $\delta-$ and Hydrogen is $\delta+$).
\end{itemize}

\subsection*{Ionic Bonds}
\begin{itemize}
    \item \textbf{Ionic Bond:} Forms when two atoms are so unequal in electronegativity that one completely \textbf{strips an electron} from the other.
    \item \textbf{Ions:} Charged atoms. \textbf{Cations} are positive (lost electrons); \textbf{Anions} are negative (gained electrons).
    \item \textbf{Salts:} Ionic compounds (e.g., NaCl). They form 3D crystal lattices, not distinct molecules. Ionic bonds are strong in dry crystals but weak in aqueous solutions.
\end{itemize}

\subsection*{Weak Chemical Interactions}
\begin{itemize}
    \item \textbf{Hydrogen Bonds:} Noncovalent attraction between a hydrogen atom ($\delta+$) covalently bonded to one electronegative atom, and a different nearby electronegative atom ($\delta-$) (usually O or N).
    \item \textbf{Van der Waals Interactions:} Weak attractions between transient positive/negative regions of molecules caused by ever-changing electron distributions (e.g., allows geckos to climb walls).
\end{itemize}

\subsection*{Molecular Shape and Function}
\begin{itemize}
    \item Shape determines biological recognition and function. 
    \item In many molecules with Carbon or Oxygen, $s$ and $p$ orbitals hybridize into a \textbf{tetrahedron} shape (e.g., $CH_4$, $H_2O$ is roughly V-shaped).
    \item \textbf{Biological Example:} Opiates (like morphine) have similar molecular shapes to natural endorphins, allowing them to bind to the same brain receptors and produce similar effects.
\end{itemize}

\section*{2.4 Chemical Reactions}
\begin{itemize}
    \item \textbf{Chemical Reactions:} The making and breaking of chemical bonds, rearranging matter. (Reactants $\rightarrow$ Products). Matter is conserved, never created nor destroyed.
    \item \textbf{Photosynthesis:} $6CO_2 + 6H_2O \rightarrow C_6H_{12}O_6 + 6O_2$. Sunlight provides the energy.
    \item \textbf{Chemical Equilibrium:} The point at which the forward and reverse reactions occur at the \textbf{same rate}. Relative concentrations of reactants and products stabilize (they do not have to be equal).
\end{itemize}

\end{document}
